% !TEX root = ../main.tex
% 本文件由 build-tex.sh 从 34-DSP指南-术语表中译.md 生成，改内容请改 md 后重新生成

\chapter{术语表}\label{ch-34}

\section{数字与符号}

\textbf{1/f noise} 1/f
噪声：一种随机噪声，频率越低幅度越大。在物理系统中普遍可见，但人们对它的了解并不深。对比参见白噪声(white
noise)。（见\hyperref[ch-09]{第 9 章}）

\textbf{−3dB cutoff frequency} −3 dB
截止频率：滤波器通带与过渡带的分界。定义为频率响应降到 −3 dB（幅度
0.707）处的频率。（见\hyperref[ch-14]{第 14 章}）

\section{A}

\textbf{``A'' law} A
律：欧洲使用的压扩标准。让量化级不等间距，使数字语音信号只用 8 位而不是
12 位就能表示。对比参见 μ 律(mu law)。（见\hyperref[ch-22]{第 22 章}）

\textbf{AC} 交流：Alternating Current
的缩写。电学术语，指信号中围绕平均值（直流）起伏的那部分。（见\hyperref[ch-02]{第
2 章}）

\textbf{Accuracy}
准确度：测量（或预测）中每次试验都会重复出现的误差。准确度受系统性（可重复的）误差限制。对比参见精密度(precision)。（见\hyperref[ch-02]{第
2 章}）

\textbf{Additivity} 可加性：线性系统必备的一条数学性质。若输入 a
产生输出 p，输入 b 产生输出 q，则输入 a+b 产生输出
p+q。（见\hyperref[ch-05]{第 5 章}）

\textbf{Aliasing}
混叠：由于采样或其他非线性操作，正弦波从一个频率变成另一个频率的过程。通常会导致信号信息丢失。（见\hyperref[ch-03]{第
3 章}）

\textbf{Amplitude modulation}
幅度调制（调幅）：无线电通信中把携带信息的信号（如音频）与载波组合起来的方法。通常把两个信号相乘来实现。（见\hyperref[ch-10]{第
10 章}）

\textbf{Analysis}
分析：傅里叶正变换，即由时域计算频域。对比参见合成(synthesis)。（见\hyperref[ch-08]{第
8 章}）

\textbf{Antialias filter}
抗混叠滤波器：放在模数转换器之前的模拟低通滤波器。去掉高于采样率一半、在转换时会发生混叠的频率。（见\hyperref[ch-03]{第
3 章}）

\textbf{ASCII}：用二进制表示字母和数字的一种方法。每个字符被分配一个
0～127 之间的数。在计算机和通信中使用极广。（见\hyperref[ch-27]{第 27
章}）

\textbf{Aspect ratio} 宽高比：图像宽度与高度之比。标准电视的宽高比为
4:3，电影为 16:9。（见\hyperref[ch-23]{第 23 章}）

\textbf{Assembly}
汇编：直接操作微处理器寄存器和内部硬件的低级编程语言。对比参见高级语言(high-level
language)。（见\hyperref[ch-04]{第 4 章}）

\textbf{Associative property of convolution} 卷积的结合律：写作
\((a[n] * b[n]) * c[n] = a[n] * (b[n] * c[n])\)。它在信号处理中很重要，因为它描述了级联各级的行为。（见\hyperref[ch-07]{第
7 章}）

\textbf{Autocorrelation}
自相关：信号与自身做相关。它有用是因为自相关的傅里叶变换就是原信号的功率谱。（见\hyperref[ch-07]{第
7 章}）

\section{B}

\textbf{Backprojection}
反投影：计算机断层成像中由各个视图重建图像的一种技术。除非与更高级的方法配合使用，否则得到的图像质量很差。（见\hyperref[ch-25]{第
25 章}）

\textbf{BASIC}：一种高级编程语言，以简单著称，也以缺点众多著称。本书的程序大多用
BASIC 写成。（见\hyperref[ch-04]{第 4 章}）

\textbf{Basilar membrane}
基底膜：耳朵里起频谱分析仪作用的小器官。它让耳蜗神经中不同的纤维受不同频率的刺激。（见\hyperref[ch-22]{第
22 章}）

\textbf{Basis functions}
基函数：一种分解所用的那组波形。例如，傅里叶分解的基函数是幅度为一的正弦波和余弦波。（见\hyperref[ch-08]{第
8 章}）

\textbf{Bessel filter}
贝塞尔滤波器：为线性相位而优化的模拟滤波器。它的阶跃响应几乎没有过冲，上升沿和下降沿相似。用来平滑时域编码的信号。（见\hyperref[ch-03]{第
3 章}）

\textbf{Bidirectional filtering}
双向滤波：用来得到零相位滤波器的递归方法。信号先从左往右滤波，再把中间结果从右往左滤波。（见\hyperref[ch-19]{第
19 章}）

\textbf{Bilinear transform} 双线性变换：把 s 平面映射到 z
平面的技术。可以把模拟滤波器转换成等效的数字滤波器。（见\hyperref[ch-33]{第
33 章}）

\textbf{Binning}
分箱：数据（或信号）的量化级非常多（例如浮点数）时，构造直方图的方法。（见\hyperref[ch-02]{第
2 章}）

\textbf{Biquad}
双二阶节：有两个极点和至多两个零点的模拟或数字系统。常级联起来构成更复杂的滤波器。（见\hyperref[ch-33]{第
33 章}）

\textbf{Bit reversal sorting} 位反转排序：FFT
中用来实现信号交错分解的算法。做法是把二进制计数的各位左右翻转。（见\hyperref[ch-12]{第
12 章}）

\textbf{Blackman window}
布莱克曼窗：用于滤波器设计和频谱分析的一条平滑曲线，按
\(0.42 - 0.5\cos(2\pi n/M) + 0.08\cos(4\pi n/M)\) 计算，n 从 0 到
M。（见\hyperref[ch-16]{第 16 章}）

\textbf{Brightness}
亮度：图像整体的明暗程度。对比参见对比度(contrast)。（见\hyperref[ch-23]{第
23 章}）

\textbf{Butterfly} 蝶形运算：FFT
中的基本计算，把两个复数变成另外两个复数。（见\hyperref[ch-12]{第 12
章}）

\textbf{Butterworth filter}
巴特沃斯滤波器：把一个频带与另一个频带分开；在保持通带平坦的前提下滚降最快；可以是模拟的，也可以是数字的。也叫最平坦(maximally
flat)滤波器。（见\hyperref[ch-03]{第 3 章}）

\section{C}

\textbf{C}：科学、工程和 DSP 中常用的编程语言。还有更高级的
C++。（见\hyperref[ch-04]{第 4 章}）

\textbf{Carrier wave}
载波：无线电信号幅度调制中的术语。指与较低频率的信息信号组合的高频正弦波。（见\hyperref[ch-10]{第
10 章}）

\textbf{Cascade}
级联：两级或更多级的组合，其中一级的输出成为下一级的输入。（见\hyperref[ch-07]{第
7 章}）

\textbf{Causal signal}
因果信号：所有负序号采样都为零的信号。（见\hyperref[ch-07]{第 7 章}）

\textbf{Causal system}
因果系统：在输入出现非零值之前输出一直为零的系统（即输入导致输出）。因果系统的冲激响应是因果信号。（见\hyperref[ch-07]{第
7 章}）

\textbf{Central Limit Theorem}
中心极限定理：统计学中的重要定理。它的一种表述是：许多随机数之和服从高斯概率分布，不论每个随机数本身服从什么分布。（见\hyperref[ch-02]{第
2 章}）

\textbf{Cepstrum} 倒谱：把
``spectrum''（频谱）一词的字母重新排列而成。用于同态处理，描述时域和频域互换后的频谱。（见\hyperref[ch-22]{第
22 章}）

\textbf{Charge coupled device (CCD)}
电荷耦合器件：电子相机中的光传感器。由一薄片硅构成，上面是一个由感光区域组成的二维阵列，这些区域叫做阱(well)。（见\hyperref[ch-23]{第
23 章}）

\textbf{Chebyshev filter}
切比雪夫滤波器：用于把一个频带与另一个频带分开。通过允许通带纹波，实现比巴特沃斯更快的滚降。可以是模拟的，也可以是数字的。（见\hyperref[ch-20]{第
20 章}）

\textbf{Chirp system}
啁啾系统：用于雷达和声呐。发射前把一个冲激转换成持续时间更长的信号，接收后再压缩回冲激。（见\hyperref[ch-11]{第
11 章}）

\textbf{Circular buffer}
循环缓冲：实时处理中使用的数据存储方法；每个新采集的采样替换存储器中最旧的采样。（见\hyperref[ch-28]{第
28 章}）

\textbf{Circular convolution}
循环卷积：频域信号相乘时可能在时域发生的混叠。时域中的每个周期都溢入相邻的周期。（见\hyperref[ch-09]{第
9 章}）

\textbf{Circularity}
循环性：信号的末尾看起来与开头相连。只考察周期信号的单个周期时就会出现这种现象。

\textbf{Classifiers}
分类量：从较大的数据集中提取出来、代表这个数据集的参数。例如：区域的大小、峰的幅度、边缘的锐度等。用于模式识别。（见\hyperref[ch-26]{第
26 章}）

\textbf{Closing}
闭运算：一种形态学运算，定义为先腐蚀、后膨胀。（见\hyperref[ch-25]{第 25
章}）

\begin{quote}
\textbf{译注：} 原书这里的定义与第 25 章正文相反。按通行定义（以及原书第
25
章），闭运算是先膨胀、后腐蚀；开运算(opening)是先腐蚀、后膨胀。原书术语表把两者的顺序写反了，见下文「Opening」条。
\end{quote}

\textbf{Cochlea} 耳蜗：耳朵中把声音转换成神经信号的器官。（原书 ``sound
in converted'' 中的 in 应为 is。）（见\hyperref[ch-22]{第 22 章}）

\textbf{Cochlear nerve}
耳蜗神经：把音频信息从耳朵传到大脑的神经。（见\hyperref[ch-22]{第 22
章}）

\textbf{Coefficient-of-variation (CV)}
变异系数：表述数据中变动（噪声）的常用方式。定义为 100\% × 标准差 /
均值。（见\hyperref[ch-02]{第 2 章}）

\textbf{Commutative property of convolution} 卷积的交换律：写作
\(a[n] * b[n] = b[n] * a[n]\)。（见\hyperref[ch-07]{第 7 章}）

\textbf{Companding} 压扩：一种 ``S'' 形非线性，使语音信号只用 8 位而不是
12 位就能数字化。欧洲用 A 律，美国用 μ 律。（见\hyperref[ch-22]{第 22
章}）

\textbf{Complex conjugation}
取共轭复数：改变复数虚部的符号。常在变量旁加一个星号表示。例如：若
\(A = 3+2j\)，则 \(A^* = 3-2j\)。（见\hyperref[ch-30]{第 30 章}）

\textbf{Complex DFT} 复数 DFT：使用复数的离散傅里叶变换。比实数 DFT
更复杂，也更强大。（见\hyperref[ch-31]{第 31 章}）

\textbf{Complex exponential} 复指数：形如 \(e^{a+bj}\)
的复数。它们在科学和工程中很有用，因为借助欧拉关系可以用它们表示正弦波。（见\hyperref[ch-30]{第
30 章}）

\textbf{Complex Fourier transform}
复傅里叶变换：傅里叶变换家族四个成员中任何一个用复数写出的形式。对比参见实傅里叶变换(real
Fourier transform)。（见\hyperref[ch-31]{第 31 章}）

\textbf{Complex numbers} 复数：实数（日常数学中用的数）加上虚数（含有 j
的数，其中 \(j = \sqrt{-1}\)）。例如：\(3+2j\)。（见\hyperref[ch-30]{第
30 章}）

\textbf{Complex plane} 复平面：复数的图形表示，实部在 x 轴，虚部在 y
轴。这和普通数用的数轴类似。（见\hyperref[ch-30]{第 30 章}）

\textbf{Composite video}
复合视频：含有同步脉冲、用来分隔场或帧的模拟电视信号。（见\hyperref[ch-23]{第
23 章}）

\textbf{Computed tomography (CT)} 计算机断层成像：由物体的 X
射线投影重建其内部图像的方法。在医学中应用很广；是 DSP
最早的应用之一。旧称 CAT 扫描仪(CAT scanner)。（见\hyperref[ch-25]{第 25
章}）

\textbf{Continuous signal}
连续信号：由连续（而非离散）变量构成的信号。例如：随时间变化的电压。常与模拟信号混用。（见\hyperref[ch-02]{第
2 章}）

\textbf{Contrast}
对比度：物体亮度与背景亮度之差。对比参见亮度(brightness)。（见\hyperref[ch-23]{第
23 章}）

\textbf{Converge}
收敛：迭代方法中的术语，表示正在向解推进（``算法正在收敛''），或已经得到解（``算法已经收敛''）。（见\hyperref[ch-26]{第
26 章}）

\textbf{Convolution integral}
卷积积分：定义连续系统中卷积的数学公式；与离散系统的卷积和相对应。（见\hyperref[ch-13]{第
13 章}）

\textbf{Convolution kernel}
卷积核：用卷积实现的滤波器的冲激响应。也叫滤波器核(filter
kernel)或核(kernel)。（见\hyperref[ch-06]{第 6 章}）

\textbf{Convolution sum}
卷积和：定义离散系统中卷积的数学公式。（见\hyperref[ch-06]{第 6 章}）

\textbf{Cooley and Tukey} 库利和图基：J.W. Cooley 和 J.W. Tukey，他们
1965 年发表的论文被认为把 FFT 带给了世界。（见\hyperref[ch-12]{第 12
章}）

\textbf{Correlation}
相关：与卷积做法相同的数学运算，只是不把其中一个信号左右翻转。这是在信号中检测已知波形的最优方法。（见\hyperref[ch-07]{第
7 章}）

\textbf{Cross-correlation}
互相关：一个信号与另一个信号做相关得到的信号。其中的峰表示两个原信号之间的相似性。另见自相关(autocorrelation)。（见\hyperref[ch-07]{第
7 章}）

\textbf{Cutoff frequency}
截止频率：在模拟和数字滤波器中，分隔通带与过渡带的频率。常取幅度降到
0.707（−3 dB）处测量。（见\hyperref[ch-14]{第 14 章}）

\textbf{CVSD} 连续可变斜率增量调制：Continuously Variable Slope Delta
modulation
的缩写，一种把语音信号转换成连续二进制流的技术。（见\hyperref[ch-03]{第
3 章}）

\section{D}

\textbf{DC} 直流：Direct Current
的缩写。电学术语，指信号中不随时间变化的部分，即平均值或均值。对比参见交流(AC)。（见\hyperref[ch-02]{第
2 章}）

\textbf{Decibel SPL} 分贝声压级：SPL 即 Sound Pressure
Level（声压级）。表示声波强度的对数刻度：0 dB SPL 刚能听到；60 dB SPL
是正常说话；140 dB SPL 会损伤耳朵。（见\hyperref[ch-22]{第 22 章}）

\begin{quote}
\textbf{译注：} 这里原书把 SPL 写作 Sound Pressure
Level（声压级），是通行写法；第 22 章正文写作 Sound Power
Level，两处不一致，以此处为准。
\end{quote}

\textbf{Decimation}
抽取：降低数字化信号的采样率。通常先做低通滤波，再丢弃部分采样。对比参见插值(interpolation)。（见\hyperref[ch-03]{第
3 章}）

\textbf{Decomposition}
分解：把一个信号拆成两个或更多相加分量的过程。常专指傅里叶正变换，即把信号拆成正弦波。（见\hyperref[ch-05]{第
5 章}）

\textbf{Deconvolution} 反卷积：卷积的逆运算：若
\(x[n] * h[n] = y[n]\)，只给定 h{[}n{]} 和 y{[}n{]}，求
x{[}n{]}。反卷积通常用频谱相除来实现。（见\hyperref[ch-17]{第 17 章}）

\textbf{Delta encoding}
差分编码：一类技术的统称，把数据存成相邻采样之差。用于模数转换、数据压缩以及许多其他应用。（见\hyperref[ch-27]{第
27 章}）

\textbf{Delta function} 冲激函数：归一化的冲激。离散冲激函数是除第 0
个采样值为一外、其余全为零的信号。连续冲激函数与之类似，但更抽象。（见\hyperref[ch-06]{第
6 章}）

\textbf{Delta-sigma}
增量-累加：在语音和音乐处理中流行的模数转换方法。使用很高的采样率、每个采样只用一位，然后再做抽取。（见\hyperref[ch-03]{第
3 章}）

\textbf{Dependent variable}
因变量：在信号中，因变量取决于自变量的值。例如：电压随时间变化时，时间是自变量，电压是因变量。（见\hyperref[ch-02]{第
2 章}）

\textbf{Difference equation}
差分方程：把输出信号的过去和当前采样，与输入信号的过去和当前采样联系起来的方程。也叫递归方程(recursion
equation)。（见\hyperref[ch-19]{第 19 章}）

\textbf{Dilation}
膨胀：一种形态学运算。作用于二值图像时，膨胀使物体变大，并可能把分开的物体连成一个。（见\hyperref[ch-25]{第
25 章}）

\textbf{Discrete cosine transform (DCT)}
离散余弦变换：傅里叶变换的一个亲戚。把信号分解成余弦波。用于数据压缩。（见\hyperref[ch-27]{第
27 章}）

\textbf{Discrete derivative}
离散导数：对离散信号的一种运算，相当于连续信号的导数。更好的名字是一阶差分(first
difference)。（见\hyperref[ch-07]{第 7 章}）

\textbf{Discrete Fourier transform (DFT)}
离散傅里叶变换：傅里叶变换家族中处理离散、周期时域信号的成员。（见\hyperref[ch-08]{第
8 章}）

\textbf{Discrete integral}
离散积分：对离散信号的一种运算，相当于连续信号的积分。更好的名字是累加和(running
sum)。（见\hyperref[ch-07]{第 7 章}）

\textbf{Discrete signal}
离散信号：使用量化变量的信号，例如存放在计算机中的数字化信号。（见\hyperref[ch-02]{第
2 章}）

\textbf{Discrete time Fourier transform (DTFT)}
离散时间傅里叶变换：傅里叶变换家族中处理离散、非周期时域信号的成员。（见\hyperref[ch-10]{第
10 章}）

\textbf{Dithering}
抖动：在模数转换之前给模拟信号加噪声，防止数字化后的信号''卡''在某一个值上。（见\hyperref[ch-03]{第
3 章}）

\textbf{Domain}
域：信号的自变量。例如，随时间变化的电压处在时域。其他常见的域有空间域（如图像）和频域（傅里叶变换的输出）。（见\hyperref[ch-02]{第
2 章}）

\textbf{Double precision} 双精度：用 64
位表示每个数的浮点记数标准。对比参见单精度(single
precision)。（见\hyperref[ch-04]{第 4 章}）

\textbf{DSP microprocessor} DSP
微处理器：专为快速数学运算设计的一类微处理器。通常带有流水线和/或哈佛结构。也叫
RISC。（见\hyperref[ch-28]{第 28 章}）

\textbf{Dynamic range}
动态范围：系统能处理的最大幅度除以系统固有的噪声。也用来指模数转换器的位数。还可以用于幅度以外的参数，参见频率动态范围(frequency
dynamic range)。

\section{E}

\textbf{Edge enhancement}
边缘增强：任何让边缘更明显的图像处理算法。也叫锐化(sharpening)操作。（见\hyperref[ch-24]{第
24 章}）

\textbf{Edge response}
边缘响应：在图像处理中，输入为边缘时系统的输出。边缘响应的锐度常用来衡量系统的分辨率。（见\hyperref[ch-25]{第
25 章}）

\textbf{Elliptic filter}
椭圆滤波器：用于把一个频带与另一个频带分开。通过允许通带和阻带都有纹波，实现很快的滚降。模拟和数字设计中都可使用。（见\hyperref[ch-20]{第
20 章}）

\textbf{End effects}
端点效应：滤波器核没有完全浸没在输入信号中，导致滤波后信号两端表现不佳。（见\hyperref[ch-06]{第
6 章}）

\textbf{Erosion}
腐蚀：一种形态学运算。作用于二值图像时，腐蚀使物体变小，并可能把一个物体断成两块或更多块。（见\hyperref[ch-25]{第
25 章}）

\textbf{Euler's relation}
欧拉关系：复数数学中最重要的公式，把正弦波、余弦波与复指数联系起来。（见\hyperref[ch-30]{第
30 章}）

\textbf{Even/odd decomposition}
偶/奇分解：把一个信号拆成另外两个信号的方法，一个偶对称，另一个奇对称。（见\hyperref[ch-05]{第
5 章}）

\textbf{Even order filter}
偶数阶滤波器：有偶数个极点的模拟或数字滤波器。

\section{F}

\textbf{False-negative}
假阴性：目标检测试验四种可能结果之一。目标存在，却被误判为不存在。（见\hyperref[ch-26]{第
26 章}）

\textbf{False-positive}
假阳性：目标检测试验四种可能结果之一。目标不存在，却被误判为存在。（见\hyperref[ch-26]{第
26 章}）

\textbf{Fast Fourier transform (FFT)}
快速傅里叶变换：计算离散傅里叶变换（DFT）的高效算法。有些情况下能把执行时间缩短几百倍。（见\hyperref[ch-12]{第
12 章}）

\textbf{FFT convolution} FFT
卷积：通过频谱相乘来卷积信号的方法。之所以叫这个名字，是因为用 FFT
在时域和频域之间高效地转换。（见\hyperref[ch-18]{第 18 章}）

\textbf{Field}
场：隔行电视先显示每帧（图像）的偶数行，再显示奇数行。偶数行叫偶场，奇数行叫奇场。（见\hyperref[ch-23]{第
23 章}）

\textbf{Filter kernel}
滤波器核：用卷积实现的滤波器的冲激响应。也叫卷积核(convolution
kernel)或核(kernel)。（见\hyperref[ch-06]{第 6 章}）

\textbf{Filtered backprojection}
滤波反投影：计算机断层成像中由各个视图重建图像的一种技术。先对各视图滤波，再做反投影。（见\hyperref[ch-25]{第
25 章}）

\textbf{Finite impulse response (FIR)}
有限冲激响应：非零值个数有限的冲激响应。常用来表示滤波器是用卷积而不是递归实现的。（见\hyperref[ch-14]{第
14 章}）

\textbf{First difference}
一阶差分：对离散信号的一种运算，模仿连续信号的一阶导数；也叫离散导数(discrete
derivative)。（见\hyperref[ch-07]{第 7 章}）

\textbf{Fixed point}
定点：计算机存储数的两种常用方式之一；通常用来存整数。对比参见浮点(floating
point)。（见\hyperref[ch-04]{第 4 章}）

\textbf{Flat-top window}
平顶窗：用于频谱分析的一种窗；能准确测量各频谱分量的幅度。可以用加窗
sinc 滤波器核来做。（见\hyperref[ch-09]{第 9 章}）

\textbf{Floating point}
浮点：计算机存储数的两种常用方式之一。浮点用一种科学记数法，把尾数乘以某个指数次幂。对比参见定点(fixed
point)。（见\hyperref[ch-04]{第 4 章}）

\textbf{Forward transform}
正变换：傅里叶变换的分析方程，由时域计算频域。对比参见逆变换(inverse
transform)。（见\hyperref[ch-08]{第 8 章}）

\textbf{Fourier reconstruction}
傅里叶重建：计算机断层成像中由各个视图计算图像的方法之一。（见\hyperref[ch-25]{第
25 章}）

\textbf{Fourier series}
傅里叶级数：傅里叶变换家族中处理连续、周期时域信号的成员。（见\hyperref[ch-13]{第
13 章}）

\textbf{Fourier transform}
傅里叶变换：一族基于把信号分解成正弦波的数学技术。在复数形式中，信号被分解成复指数。（见\hyperref[ch-08]{第
8 章}）

\textbf{Fourier transform pair}
傅里叶变换对：时域和频域中相互对应的波形。例如矩形脉冲和 sinc
函数。（见\hyperref[ch-11]{第 11 章}）

\textbf{Fovea}
中央凹：眼睛视网膜上为高分辨率视觉而优化的一小块区域。（见\hyperref[ch-23]{第
23 章}）

\textbf{Frame} 帧：电视信号中的一幅单独图像。NTSC 电视标准每秒 30
帧。（见\hyperref[ch-23]{第 23 章}）

\textbf{Frame grabber}
帧采集卡：用来把电视信号中的一帧（图像）数字化并存储下来的模数转换器。（见\hyperref[ch-23]{第
23 章}）

\textbf{Frequency domain}
频域：以频率为自变量的信号。傅里叶变换的输出。（见\hyperref[ch-08]{第 8
章}）

\textbf{Frequency domain aliasing}
频域混叠：由于在时域中进行某种操作而在频域发生的混叠。采样时的混叠就是一个例子。（见\hyperref[ch-10]{第
10 章}）

\textbf{Frequency domain convolution}
频域卷积：通过信号频谱相乘来完成的卷积。（见\hyperref[ch-09]{第 9 章}）

\textbf{Frequency domain encoding}
频域编码：信息在信号中编码的两种主要方式之一。信息包含在信号各正弦分量的幅度、频率和相位中。音频信号是最好的例子。对比参见时域编码(time
domain encoding)。（见\hyperref[ch-14]{第 14 章}）

\textbf{Frequency domain multiplexing}
频分复用：把信号移到频谱的不同位置、以便同时传输的组合方法。（见\hyperref[ch-10]{第
10 章}）

\textbf{Frequency dynamic range}
频率动态范围：系统能处理的最高频率与最低频率之比。模拟系统的频率动态范围通常比数字系统大得多。（见\hyperref[ch-21]{第
21 章}）

\textbf{Frequency resolution}
频率分辨率：辨别或分开靠得很近的频率的能力。（见\hyperref[ch-09]{第 9
章}）

\textbf{Frequency response}
频率响应：正弦波通过线性系统时经历的幅度和相位变化。通常表示为频率的函数。常通过对冲激响应做傅里叶变换求得。（见\hyperref[ch-09]{第
9 章}）

\textbf{Fricative} 擦音：起源于气流湍动产生的随机噪声的人声，例如
s、f、sh、z、v 和 th。对比参见浊音(voiced)。（见\hyperref[ch-22]{第 22
章}）

\textbf{Full-width-at-half-maximum (FWHM)}
半高全宽：测量信号中峰宽度的常用方法。在峰最大幅度的一半处测量峰的宽度。（见\hyperref[ch-25]{第
25 章}）

\textbf{Fundamental frequency}
基频：周期波形重复自身的频率。对比参见谐波(harmonic)。（见\hyperref[ch-11]{第
11 章}）

\section{G}

\textbf{Gamma curve}
伽马曲线：把存储的像素值与它在显示图像中呈现的亮度联系起来的数学函数或查找表。也叫灰度变换(grayscale
transform)。（见\hyperref[ch-23]{第 23 章}）

\textbf{Gaussian} 高斯：钟形曲线，一般形式为
\(e^{-x^2}\)。高斯曲线有许多独特的性质。也叫正态分布(normal
distribution)。（见\hyperref[ch-02]{第 2 章}）

\begin{quote}
\textbf{译注：} 原书此处（以及下文「Normal distribution」条）印作
\(e^{x^2}\)，漏了负号；钟形曲线应为
\(e^{-x^2}\)。我放大原书页面确认指数里没有负号，这是原书印刷错误，不是文字抽取丢失。
\end{quote}

\textbf{Gibbs effect}
吉布斯效应：信号在一个域中被截断时，另一个域中的边缘和拐角处会出现振铃和过冲。（见\hyperref[ch-11]{第
11 章}）

\textbf{GIF}：一种常见的图像文件格式，使用
LZW（无损）压缩。在万维网上广泛用于图形。对比参见 TIFF 和
JPEG。（见\hyperref[ch-27]{第 27 章}）

\textbf{Grayscale image}
灰度图像：每个像素以黑白之间的灰度显示的数字图像；也叫黑白图像。（见\hyperref[ch-23]{第
23 章}）

\textbf{Grayscale stretch}
灰度拉伸：大幅提高数字图像的对比度，以便细看一小段量化级。这段范围以外的量化级显示为饱和的黑或白。（见\hyperref[ch-23]{第
23 章}）

\textbf{Grayscale transform}
灰度变换：存储的像素值与显示图像中呈现的亮度之间的转换函数。也叫伽马曲线(gamma
curve)。（见\hyperref[ch-23]{第 23 章}）

\section{H}

\textbf{Halftone}
半色调：在纸上印刷图像的常用方法。用各种小黑点图案形成灰度。彩色半色调使用红、绿、蓝点。（见\hyperref[ch-23]{第
23 章}）

\begin{quote}
\textbf{译注：}
彩色印刷实际使用的是青、品红、黄（通常再加黑）四色网点；红、绿、蓝是显示器的加色三原色。
\end{quote}

\textbf{Hamming window}
汉明窗：用于滤波器设计和频谱分析的一条平滑曲线，按
\(0.54 - 0.46\cos(2\pi n/M)\) 计算，n 从 0 到 M。（见\hyperref[ch-16]{第
16 章}）

\textbf{Harmonics}
谐波：周期信号的频率分量，总是基频的整数倍。基频是一次谐波，两倍基频是二次谐波，依此类推。（见\hyperref[ch-11]{第
11 章}）

\textbf{Harvard Architecture}
哈佛结构：程序和数据放在不同存储器、通过不同总线访问的计算机内部布局；在用于
DSP 的微处理器中很常见。对比参见冯·诺依曼结构(Von Neumann
Architecture)。（见\hyperref[ch-28]{第 28 章}）

\textbf{High fidelity} 高保真：高质量的音乐重放，例如 CD
播放机提供的效果。（见\hyperref[ch-22]{第 22 章}）

\textbf{High-level language} 高级语言：C、BASIC、FORTRAN
之类的编程语言。（见\hyperref[ch-04]{第 4 章}）

\textbf{High-speed convolution} 高速卷积：FFT
卷积的另一个名字。（见\hyperref[ch-18]{第 18 章}）

\textbf{Hilbert transformer}
希尔伯特变换器：对所有频率，频率响应都是幅度 = 1、相位 = 90°
的系统。在通信系统中用于调制。可以是模拟的，也可以是数字的。（见\hyperref[ch-33]{第
33 章}）

\textbf{Histogram equalization}
直方图均衡化：用图像的积分直方图作为灰度变换来处理图像。原理是让图像中面积大的部分比面积小的部分获得更高的对比度。（见\hyperref[ch-23]{第
23 章}）

\textbf{Histogram} 直方图：显示信号中各个值的分布。x
轴是采样可能取的值，y 轴是取每个值的采样个数。（见\hyperref[ch-02]{第 2
章}）

\textbf{Homogeneity} 齐次性：所有线性系统都具有的一条数学性质。若输入
x{[}n{]} 产生输出 y{[}n{]}，则对任意常数 k，输入 k·x{[}n{]} 产生输出
k·y{[}n{]}。（见\hyperref[ch-05]{第 5 章}）

\textbf{Homomorphic}
同态：分离以非线性方式（例如相乘或卷积）组合在一起的信号的 DSP
技术。通过适当的变换把非线性问题转化成线性问题。（见\hyperref[ch-22]{第
22 章}）

\textbf{Huffman encoding}
霍夫曼编码：给常见字符分配较少位、给少见字符分配较多位的数据压缩方法。（见\hyperref[ch-27]{第
27 章}）

\textbf{Hyperspace}
超空间：目标检测和神经网络分析中的术语。一个参数可以用图形表示为一条线，两个参数是一个平面，三个参数是一个空间，三个以上的参数是一个超空间。（见\hyperref[ch-26]{第
26 章}）

\section{I}

\textbf{Imaginary part} 虚部：复数中带 j 的那部分，例如 \(3+2j\) 中的
2。在实傅里叶变换中，虚部也指频域中存放正弦波幅度的那部分，尽管那里并不用
j。（见\hyperref[ch-08]{第 8 章}）

\textbf{Impulse}
冲激：除一个非常短的脉冲外全为零的信号。对离散信号，这个脉冲就是一个非零采样。对连续信号，脉冲宽度必须远小于所用系统的固有响应。（见\hyperref[ch-05]{第
5 章}）

\textbf{Impulse decomposition} 冲激分解：把一个 N 点信号拆成 N
个信号，每个信号只含原信号的一个采样，其余采样都为零。这是卷积的基础。（见\hyperref[ch-05]{第
5 章}）

\textbf{Impulse response}
冲激响应：输入为归一化冲激（冲激函数）时系统的输出。（见\hyperref[ch-06]{第
6 章}）

\textbf{Impulse train}
冲激串：由一系列等间距冲激组成的信号。（见\hyperref[ch-03]{第 3 章}）

\textbf{Independent variable}
自变量：在信号中，因变量取决于自变量的值。例如：电压随时间变化时，时间是自变量，电压是因变量。（见\hyperref[ch-02]{第
2 章}）

\textbf{Infinite impulse response (IIR)}
无限冲激响应：非零值个数无限的冲激响应，例如衰减的指数。常用来表示滤波器是用递归而不是卷积实现的。（见\hyperref[ch-19]{第
19 章}）

\textbf{Integers} 整数：\ldots，−2，−1，0，1，2，\ldots{}
这样的数。也指以定点格式存储的数。对比参见浮点(floating
point)。（见\hyperref[ch-04]{第 4 章}）

\textbf{Interlaced decomposition}
交错分解：把信号拆成偶数号采样和奇数号采样。用于
FFT。（见\hyperref[ch-12]{第 12 章}）

\textbf{Interlaced video}
隔行视频：先显示每幅图像的偶数行、再显示奇数行的视频信号。用于电视，是为减少闪烁而发展出来的。（见\hyperref[ch-23]{第
23 章}）

\textbf{Interpolation}
插值：提高数字化信号的采样率。通常在原采样之间插入零，再用低通滤波器滤波。对比参见抽取(decimation)。（见\hyperref[ch-03]{第
3 章}）

\textbf{Inverse transform}
逆变换：傅里叶变换的合成方程，由频域计算时域。对比参见正变换(forward
transform)。（见\hyperref[ch-08]{第 8 章}）

\textbf{Iterative}
迭代：沿正确方向逐步调整变量、直到收敛的求解方法。用于 CT
重建和神经网络。（见\hyperref[ch-26]{第 26 章}）

\section{J}

\textbf{JPEG}：一种常见的图像文件格式，使用变换（有损）压缩。在万维网上广泛用于图形。对比参见
GIF 和 TIFF。（见\hyperref[ch-27]{第 27 章}）

\section{K}

\textbf{Kernel} 核：用卷积实现的滤波器的冲激响应。也叫卷积核(convolution
kernel)或滤波器核(filter kernel)。（见\hyperref[ch-06]{第 6 章}）

\section{L}

\textbf{Laplace transform}
拉普拉斯变换：分析由微分方程支配的系统的数学方法。是设计电路（例如模拟滤波器）的主要工具。把时域信号变成
s 域信号。（见\hyperref[ch-32]{第 32 章}）

\textbf{Learning algorithm}
学习算法：根据网络应当如何工作的样例，求出一组神经网络权重的过程。（见\hyperref[ch-26]{第
26 章}）

\textbf{Line pair} 线对：成像领域对''周期''的说法。例如，每毫米 5
个周期等于每毫米 5 个线对。（见\hyperref[ch-25]{第 25 章}）

\textbf{Line pair gauge}
线对卡：测量成像系统分辨率的器件。上面有一组明暗相间的线条，在一端越来越密。（见\hyperref[ch-25]{第
25 章}）

\textbf{Line spread function (LSF)}
线扩散函数：成像系统对输入图像中一条细线的响应。（见\hyperref[ch-25]{第
25 章}）

\textbf{Linear phase}
线性相位：相位是一条直线的系统。它通常很重要，因为这意味着冲激响应左右对称，使输出信号的上升沿和下降沿看起来一样。另见零相位(zero
phase)。（见\hyperref[ch-07]{第 7 章}）

\textbf{Linear system}
线性系统：按定义，具有可加性和齐次性的系统。（见\hyperref[ch-05]{第 5
章}）

\textbf{Lossless compression}
无损压缩：能精确还原原始数据的数据压缩技术，例如 LZW
压缩。（见\hyperref[ch-27]{第 27 章}）

\textbf{Lossy compression}
有损压缩：只能还原原始数据近似值的数据压缩方法。这样能得到更高的压缩比。JPEG
就是一例。（见\hyperref[ch-27]{第 27 章}）

\section{M}

\textbf{Matched filtering}
匹配滤波：确定已知图案在信号中出现在哪里（或是否出现）的方法。匹配滤波基于相关，但用卷积实现。（原书
``a know pattern'' 中的 know 应为 known。）（见\hyperref[ch-07]{第 7
章}）

\textbf{Mathematical equivalence}
数学等价：用复数表示实际问题的一种方式。基于把正弦波与复指数等同起来的欧拉关系。对比参见代换法(substitution)。（见\hyperref[ch-30]{第
30 章}）

\textbf{Mean} 均值：信号或其他一组数据的平均值。（见\hyperref[ch-02]{第
2 章}）

\textbf{Memoryless}
无记忆：输出的当前值只取决于输入的当前值、与过去值无关的系统。（见\hyperref[ch-05]{第
5 章}）

\textbf{MFLOPS}：Million-Floating-Point-Operations-Per-Second（每秒百万次浮点运算）的缩写；表示计算机速度的常用方式。对比参见
MIPS。

\textbf{MIPS}：Million-Instructions-Per-Second（每秒百万条指令）的缩写；表示计算机速度的常用方式。对比参见
MFLOPS。

\textbf{Mixed signal}
混合信号：同时包含模拟和数字电路的集成电路，例如集成在数字信号处理器上的模数转换器。（见\hyperref[ch-28]{第
28 章}）

\textbf{Modulation transfer function (MTF)}
调制传递函数：成像领域对频率响应的说法。（见\hyperref[ch-25]{第 25 章}）

\textbf{Morphing}
变形动画：把一幅图像逐渐变形成另一种形态。用于特效，例如人变成狼人。（见\hyperref[ch-23]{第
23 章}）

\textbf{Morphological}
形态学的：通常指对二值图像做的简单非线性运算，例如腐蚀和膨胀。（见\hyperref[ch-25]{第
25 章}）

\textbf{Moving average filter}
移动平均滤波器：输出信号的每个采样是输入信号中许多相邻采样的平均值。可以用卷积或递归实现。（见\hyperref[ch-15]{第
15 章}）

\textbf{MPEG}：用于视频（例如数字电视）的压缩标准。（见\hyperref[ch-27]{第
27 章}）

\textbf{Mu law} μ
律：美国使用的压扩标准。让量化级不等间距，使数字语音信号只用 8 位而不是
12 位就能表示。对比参见 A 律(``A'' law)。（见\hyperref[ch-22]{第 22
章}）

\textbf{Multiplexing}
多路复用：把两个或更多信号组合在一起传输。实现方式有很多种。（原书 ``two
or move'' 中的 move 应为 more。）

\textbf{Multirate}
多速率：使用不止一个采样率的系统。常用于模数和数模转换，用更少的电路获得更好的性能。（见\hyperref[ch-03]{第
3 章}）

\section{N}

\textbf{Natural frequency}
自然频率：以弧度每秒（而不是周期每秒，即赫兹）表示的频率。把频率（赫兹）乘以
2π 就得到自然频率。（见\hyperref[ch-08]{第 8 章}）

\textbf{Negative frequencies} 负频率：正弦波既可以写成正频率
\(\cos(\omega t)\)，也可以写成负频率
\(\cos(-\omega t)\)。复傅里叶变换包含负频率，因而更强大。（见\hyperref[ch-10]{第
10 章}）

\textbf{Normal distribution} 正态分布：形如 \(e^{-x^2}\)
的钟形曲线。也叫高斯分布(Gaussian)。（原书同样漏印负号，见「Gaussian」条译注。）（见\hyperref[ch-02]{第
2 章}）

\textbf{NTSC}：美国、日本等国使用的电视标准。对比参见 PAL 和
SECAM。（见\hyperref[ch-23]{第 23 章}）

\textbf{Nyquist frequency, Nyquist rate}
奈奎斯特频率、奈奎斯特速率：这两个术语与采样定理有关，但不同作者用法不同。它们可以指四种不同的东西：信号所含的最高频率、这个频率的两倍、采样率，或采样率的一半。（见\hyperref[ch-03]{第
3 章}）

\section{O}

\textbf{Octave} 八度：频率相差两倍。（见\hyperref[ch-22]{第 22 章}）

\textbf{Odd order filter} 奇数阶滤波器：有奇数个极点的模拟或数字滤波器。

\textbf{Opening}
开运算：一种形态学运算，定义为先膨胀、后腐蚀。（见\hyperref[ch-25]{第 25
章}）

\begin{quote}
\textbf{译注：}
与「Closing」条相同，原书术语表把顺序写反了。通行定义（以及原书第 25
章）是：开运算先腐蚀、后膨胀；闭运算先膨胀、后腐蚀。
\end{quote}

\textbf{Optimal filter}
最优滤波器：在某个特定意义上''最好''的滤波器。例如，维纳滤波器能得到最优信噪比，匹配滤波器对目标检测最优。（见\hyperref[ch-17]{第
17 章}）

\textbf{Overlap add}
重叠相加：把长信号拆成若干段来处理的方法。（见\hyperref[ch-18]{第 18
章}）

\section{P}

\textbf{PAL}：欧洲使用的电视标准。对比参见 NTSC。（见\hyperref[ch-23]{第
23 章}）

\textbf{Parallel stages}
并联级：两级或更多级的组合，它们输入相同、输出相加。（见\hyperref[ch-07]{第
7 章}）

\textbf{Parameter space}
参数空间：目标检测的行话。一个参数可以用图形表示为一条线，两个参数是一个平面，三个参数是一个空间，三个以上的参数是一个超空间。（见\hyperref[ch-26]{第
26 章}）

\textbf{Parseval's relation}
帕塞瓦尔关系：把时域能量与频域能量联系起来的公式。（见\hyperref[ch-10]{第
10 章}）

\textbf{Passband}
通带：滤波器设计成让其不受改变地通过的那段频率。（见\hyperref[ch-14]{第
14 章}）

\textbf{Passive sonar}
被动声呐：依靠潜艇和其他水下物体自身发出的声音来探测它们。用于隐蔽监视。

\textbf{Phasor transform} 相量变换：用复数求 RLC
电路频率响应的方法。电阻、电容和电感分别变成 R、\(-j/\omega C\) 和
\(j\omega L\)。（见\hyperref[ch-30]{第 30 章}）

\textbf{Pillbox}
圆柱形核：图像处理中滤波器核的一种形状：一个值为常数的圆形区域，周围是零。（见\hyperref[ch-24]{第
24 章}）

\textbf{Pitch}
音高：人对连续音调基频的感知。对比参见音色(timbre)。（见\hyperref[ch-22]{第
22 章}）

\textbf{Pixel} 像素：``picture
element''（图像元素）的缩写。数字图像中的单个采样。（见\hyperref[ch-23]{第
23 章}）

\textbf{Point spread function (PSF)}
点扩散函数：成像领域对冲激响应的说法。（见\hyperref[ch-24]{第 24 章}）

\textbf{Pointer} 指针：值为另一个变量地址的变量。（见\hyperref[ch-28]{第
28 章}）

\textbf{Poisson statistics} 泊松统计：信号由有限个粒子（例如 X
射线、光子或电子）表示时，其值产生的波动。也叫泊松噪声(Poisson
noise)和统计噪声(statistical noise)。（见\hyperref[ch-25]{第 25 章}）

\textbf{Polar form}
极坐标形式：用幅度和相位表示正弦波：\(M\cos(\omega t + \phi)\)，其中 M
是幅度，φ 是相位。对比参见直角坐标形式(rectangular
form)。（见\hyperref[ch-08]{第 8 章}）

\textbf{Pole} 极点：拉普拉斯变换和 z 变换中的术语。把 s 域或 z
域的传递函数写成一个多项式除以另一个多项式时，分母的根是系统的极点，分子的根是零点。（见\hyperref[ch-32]{第
32 章}）

\textbf{Pole-zero diagram} 零极点图：拉普拉斯变换和 z 变换中的术语。在 s
平面或 z 平面上标出极点和零点位置的图。（见\hyperref[ch-32]{第 32 章}）

\textbf{Precision}
精密度：测量或预测中每次试验都不同、不可重复的误差。精密度由随机误差决定。对比参见准确度(accuracy)。（见\hyperref[ch-02]{第
2 章}）

\textbf{Probability distribution function (pdf)}
概率分布函数：给出连续变量取某个值的概率。（见\hyperref[ch-02]{第 2
章}）

\textbf{Probability mass function (pmf)}
概率质量函数：给出离散变量取某个值的概率。对比参见概率分布函数(pdf)。（见\hyperref[ch-02]{第
2 章}）

\textbf{Pulse response}
脉冲响应：输入为脉冲时系统的输出。（见\hyperref[ch-19]{第 19 章}）

\section{Q}

\textbf{Quantization error}
量化误差：信号量化时引入的误差。大多数情况下，最大误差为 ±½
LSB，均方根误差为 \(1/\sqrt{12}\) LSB。也叫量化噪声(quantization
noise)。（见\hyperref[ch-03]{第 3 章}）

\section{R}

\textbf{Random error}
随机误差：测量或预测中每次试验都不同、不可重复的误差。它决定精密度。对比参见系统误差(systematic
error)。（见\hyperref[ch-02]{第 2 章}）

\textbf{Radar} 雷达：Radio Detection And
Ranging（无线电探测与测距）的缩写。用无线电波探测飞机的回声定位技术。

\textbf{Real DFT} 实数
DFT：只使用实数（普通数）的离散傅里叶变换。不如复数 DFT
强大，但更简单。对比参见复数 DFT(complex DFT)。（见\hyperref[ch-08]{第 8
章}）

\textbf{Real FFT} 实数 FFT：FFT
的一种改进版本。时域完全为实数（即时域虚部为零）时，比标准 FFT 快约
30\%。（第 12 章译者在本机实测，实数 FFT
约快一倍。）（见\hyperref[ch-12]{第 12 章}）

\textbf{Real Fourier transform}
实傅里叶变换：傅里叶变换家族中只使用实数（而不是虚数或复数）的任何成员。对比参见复傅里叶变换(complex
Fourier transform)。（见\hyperref[ch-31]{第 31 章}）

\textbf{Real part} 实部：复数中不带 j 的那部分，例如 \(3+2j\) 中的
3。在实傅里叶变换中，实部指频域中存放余弦波幅度的那部分，尽管那里并没有
j。（见\hyperref[ch-08]{第 8 章}）

\textbf{Real time processing}
实时处理：数据一采集到就处理，而不是存起来以后再用。例如：控制长途电话回声的
DSP 算法。

\textbf{Reconstruction filter}
重建滤波器：放在数模转换器之后的模拟低通滤波器。去掉高于采样率一半的频率，使阶梯状的波形变平滑。（见\hyperref[ch-03]{第
3 章}）

\textbf{Rectangular form} 直角坐标形式：把正弦波表示为
\(A\cos(\omega t) + B\sin(\omega t)\)，其中 A 叫实部，B
叫虚部（尽管它们并不是虚数）。（见\hyperref[ch-08]{第 8 章}）

\textbf{Rectangular window}
矩形窗：一组相邻点的值为一、其余为零的信号。通常与另一个信号相乘，以选出要处理的那一段。（见\hyperref[ch-09]{第
9 章}）

\textbf{Recursion coefficients}
递归系数：递归方程中使用的加权值。递归系数决定递归（IIR）滤波器的特性。（见\hyperref[ch-19]{第
19 章}）

\textbf{Recursion equation}
递归方程：把输出信号的过去和当前采样，与输入信号的过去和当前值联系起来的方程。也叫差分方程(difference
equation)。（见\hyperref[ch-19]{第 19 章}）

\textbf{Region-of-convergence} 收敛域：拉普拉斯变换和 z 变换中的术语。指
s 平面和 z 平面上取值有定义的那些区域。（见\hyperref[ch-32]{第 32 章}）

\textbf{RGB encoding} RGB
编码：给每个像素指定红、绿、蓝的分量来表示彩色图像。（见\hyperref[ch-23]{第
23 章}）

\textbf{RISC} 精简指令集计算机：Reduced Instruction Set Computer
的缩写，也叫 DSP
微处理器。编程指令较少，使数学运算速度高得多。与之相反的是复杂指令集计算机(Complex
Instruction Set Computer)，例如奔腾。

\textbf{ROC curve} ROC
曲线：显示阈值选择如何影响目标检测性能的图。（见\hyperref[ch-26]{第 26
章}）

\textbf{Roll-off}
滚降：描述滤波器通带与阻带之间过渡陡峭程度的行话。滚降快表示过渡陡；滚降慢表示过渡平缓。（见\hyperref[ch-14]{第
14 章}）

\textbf{Root-mean-square (rms)}
均方根：用来表示信号围绕零的起伏。在电子学中常用。定义为各平方值的平均的平方根。对比参见标准差(standard
deviation)。（见\hyperref[ch-02]{第 2 章}）

\textbf{Round-off noise}
舍入噪声：把数学运算结果舍入到最近量化级引起的误差。

\textbf{Row major order}
行优先顺序：把图像转换成串行形式的一种顺序。和书写英文一样：第一行从左到右，第二行从左到右，依此类推。（见\hyperref[ch-23]{第
23 章}）

\textbf{Run-length encoding}
游程编码：一种简单的数据压缩技术，有许多变体。连续重复多次的字符被替换成表示该字符和游程长度的代码。（见\hyperref[ch-27]{第
27 章}）

\textbf{Running sum}
累加和：对离散信号的一种运算，模仿连续信号的积分。也叫离散积分(discrete
integral)。（见\hyperref[ch-07]{第 7 章}）

\section{S}

\textbf{s-domain} s 域：由拉普拉斯变换定义的域。也叫 s
平面(s-plane)。（见\hyperref[ch-32]{第 32 章}）

\textbf{Sample spacing}
采样间距：连续图像数字化时采样之间的间隔。定义为像素中心到中心的距离。（见\hyperref[ch-25]{第
25 章}）

\textbf{Sampling aperture}
采样孔径：数字化时，连续图像中对单个像素有贡献的区域。一般与采样间距差不多大。（见\hyperref[ch-25]{第
25 章}）

\textbf{Sampling theorem} 采样定理：如果一个由低于 f
的频率组成的连续信号以 2f
采样，连续信号中的全部信息都会保留在采样后的信号中。常称为香农采样定理或奈奎斯特采样定理。（见\hyperref[ch-03]{第
3 章}）

\textbf{SECAM}：欧洲使用的电视标准。对比参见
NTSC。（见\hyperref[ch-23]{第 23 章}）

\textbf{Seismology} 地震学：地球物理学中研究地球力学性质的分支。

\textbf{Separable}
可分离：能表示为其纵向剖面与横向剖面之积的图像。用于提高图像卷积的速度。（见\hyperref[ch-24]{第
24 章}）

\textbf{Sharpening}
锐化：让边缘更陡峭的图像处理操作。（见\hyperref[ch-24]{第 24 章}）

\textbf{Shift and subtract}
平移相减：产生三维或浮雕效果的图像处理操作。（见\hyperref[ch-24]{第 24
章}）

\textbf{Shift invariance}
平移不变性：许多系统具有的性质。输入信号的平移只会引起输出信号同样的平移。意味着系统的特性不随时间（或其他自变量）变化。（见\hyperref[ch-05]{第
5 章}）

\textbf{Sigmoid} sigmoid 函数：神经网络中使用的 ``S''
形曲线。（见\hyperref[ch-26]{第 26 章}）

\textbf{Signal}
信号：描述一个参数怎样随另一个参数变化。例如：随时间变化的电压。（见\hyperref[ch-02]{第
2 章}）

\textbf{Signal restoration}
信号复原：把以某种方式被改变或劣化的信号恢复成原来的样子。滤波的主要用途之一。（见\hyperref[ch-14]{第
14 章}）

\textbf{Sinc function} sinc 函数：正式定义为
\(\mathrm{sinc}(a) = \sin(\pi a)/(\pi a)\)。π
常被并入其他变量，使它写成一般形式
\(\sin(x)/x\)。它重要是因为它是矩形脉冲的傅里叶变换。（见\hyperref[ch-11]{第
11 章}）

\textbf{Single precision} 单精度：用 32
位表示每个数的浮点记数法。对比参见双精度(double
precision)。（见\hyperref[ch-04]{第 4 章}）

\textbf{Single-pole digital filters} 单极点数字滤波器：模仿电子学中 RC
高通和低通滤波器的简单递归滤波器。（见\hyperref[ch-19]{第 19 章}）

\textbf{Sinusoidal fidelity}
正弦保真性：线性系统的一条重要性质。正弦输入只能产生正弦输出；幅度和相位可能变化，但频率保持不变。（见\hyperref[ch-05]{第
5 章}）

\textbf{Sonar} 声呐：Sound Navigation And
Ranging（声音导航与测距）的缩写。用声音探测潜艇和其他水下物体。主动声呐用回声定位，被动声呐只听。

\textbf{Source code}
源代码：程序员写出的那种形式的计算机程序；有别于可在计算机上直接运行的可执行代码。（见\hyperref[ch-04]{第
4 章}）

\textbf{Spatial domain}
空间域：以距离（空间）为自变量的信号。图像就是空间域中的信号。（见\hyperref[ch-02]{第
2 章}）

\textbf{Spectral analysis}
频谱分析：通过考察信号各正弦分量的幅度、频率和相位来理解信号。频谱分析的主要工具是傅里叶变换。（见\hyperref[ch-09]{第
9 章}）

\textbf{Spectral inversion}
频谱反转：改变滤波器核、使对应的频率响应上下翻转的方法。可以把低通滤波器变成高通，把带通变成带阻，等等。（见\hyperref[ch-14]{第
14 章}）

\textbf{Spectral leakage}
频谱泄漏：频谱分析中的术语。由于只能对有限长的信号做
DFT，正弦波的频谱是一个带尾巴的峰。这些尾巴叫做主峰的泄漏。（见\hyperref[ch-09]{第
9 章}）

\textbf{Spectral reversal}
频谱翻转：改变滤波器核、使对应的频率响应左右翻转的技术。可以把低通滤波器变成高通滤波器。（见\hyperref[ch-14]{第
14 章}）

\textbf{Spectrogram}
语谱图：测量音频频谱随时间如何变化。通常显示为一幅图像。也叫声纹(voiceprint)。（见\hyperref[ch-22]{第
22 章}）

\textbf{Standard deviation}
标准差：表示信号围绕平均值起伏的一种方式。定义为各偏差平方的平均的平方根，偏差是采样值与均值之差。对比参见均方根(root-mean-square)。（见\hyperref[ch-02]{第
2 章}）

\textbf{Static linearity}
静态线性：指线性系统在信号不变化（即直流或静态）时的表现。此时输出等于输入乘以一个常数。（见\hyperref[ch-05]{第
5 章}）

\textbf{Statistical noise} 统计噪声：信号由有限个粒子（例如 X
射线、电子或光子）表示时，其值产生的波动。也叫泊松统计(Poisson
statistics)和泊松噪声(Poisson noise)。（见\hyperref[ch-25]{第 25 章}）

\textbf{Steepest descent}
最速下降：设计迭代算法的一种策略。好比总朝下坡方向走，来找到山谷的谷底。（见\hyperref[ch-26]{第
26 章}）

\textbf{Step response}
阶跃响应：输入为阶跃函数时系统的输出。（见\hyperref[ch-14]{第 14 章}）

\textbf{Stopband}
阻带：滤波器设计成要阻挡的那段频率。（见\hyperref[ch-14]{第 14 章}）

\textbf{Stopband attenuation}
阻带衰减：阻带中各频率的幅度被削减的程度，通常用分贝表示。用来描述滤波器的性能。（见\hyperref[ch-14]{第
14 章}）

\textbf{Substitution}
代换法：用复数表示物理问题（例如电路设计）的一种方式。这种方法先加入 j
项把物理问题变成复数形式，再去掉 j
项回到原来的形式。对比参见数学等价(mathematical
equivalence)。（见\hyperref[ch-30]{第 30 章}）

\textbf{Switched capacitor filter}
开关电容滤波器：用快速开关代替电阻的模拟滤波器。做成易用的集成电路。常用作模数转换器的抗混叠滤波器和数模转换器的重建滤波器。（见\hyperref[ch-03]{第
3 章}）

\textbf{Synthesis}
合成：傅里叶逆变换，由频域计算时域。对比参见分析(analysis)。（见\hyperref[ch-08]{第
8 章}）

\textbf{System}
系统：任何对输入信号作出响应、产生输出信号的过程。（见\hyperref[ch-05]{第
5 章}）

\textbf{Systematic error}
系统误差：测量或预测中每次试验都会重复出现的误差。系统误差决定准确度。对比参见随机误差(random
error)。（见\hyperref[ch-02]{第 2 章}）

\section{T}

\textbf{Target detection}
目标检测：根据测量值判断某个物体或状况是否存在。（见\hyperref[ch-26]{第
26 章}）

\textbf{TIFF}：文字处理等程序中使用的一种常见图像文件格式。通常不压缩，但可以选择
LZW 压缩。对比参见 GIF 和 JPEG。（见\hyperref[ch-27]{第 27 章}）

\textbf{Timbre}
音色：人对声音中谐波的感知。对比参见音高(pitch)。（见\hyperref[ch-22]{第
22 章}）

\textbf{Time domain}
时域：以时间为自变量的信号。也泛指采集数据时所在的任何域。（见\hyperref[ch-02]{第
2 章}）

\textbf{Time domain aliasing}
时域混叠：在频域中进行某种操作时，在时域发生的混叠。循环卷积就是一个例子。（见\hyperref[ch-10]{第
10 章}）

\textbf{Time domain encoding}
时域编码：信号信息包含在波形的形状中。对比参见频域编码(frequency domain
encoding)。（见\hyperref[ch-14]{第 14 章}）

\textbf{Transfer function}
传递函数：输出信号除以输入信号。根据信号的表示方式不同，它有几种不同形式。例如，在
s 域和 z
域中，它是一个多项式除以另一个多项式，可以用极点和零点来表示。（见\hyperref[ch-32]{第
32 章}）

\textbf{Transform} 变换：把一组数据变成另一组数据的过程、方程或算法。

\textbf{Transform compression}
变换压缩：基于给高频分配较少位的数据压缩技术。JPEG
是最好的例子。（见\hyperref[ch-27]{第 27 章}）

\textbf{Transition band}
过渡带：滤波器行话；通带与阻带之间、发生滚降的那段频率。（见\hyperref[ch-14]{第
14 章}）

\textbf{True-negative}
真阴性：目标检测试验四种可能结果之一。目标不存在，并被正确判为不存在。（见\hyperref[ch-26]{第
26 章}）

\textbf{True-positive}
真阳性：目标检测试验四种可能结果之一。目标存在，并被正确判为存在。（见\hyperref[ch-26]{第
26 章}）

\section{U}

\textbf{Unit circle} 单位圆：z 平面上 r = 1
的圆。沿这个圆的值就是系统的频率响应。（见\hyperref[ch-33]{第 33 章}）

\textbf{Unit impulse}
单位冲激：冲激函数的另一个名字。（见\hyperref[ch-06]{第 6 章}）

\section{V}

\textbf{Von Neumann Architecture}
冯·诺依曼结构：程序和数据放在同一个存储器中的计算机内部布局；非常常见。对比参见哈佛结构(Harvard
Architecture)。（见\hyperref[ch-28]{第 28 章}）

\textbf{Voiced}
浊音：起源于气流脉冲通过声带的人声。元音就是浊音的例子。对比参见擦音(fricative)。（见\hyperref[ch-22]{第
22 章}）

\section{W}

\textbf{Well} 阱：potential well（势阱）的简称；CCD
中感光的区域。（见\hyperref[ch-23]{第 23 章}）

\textbf{White noise}
白噪声：频谱平坦的随机噪声。当时域中每个采样都不包含关于其他采样的任何信息时就会出现。对比参见
1/f 噪声(1/f noise)。（见\hyperref[ch-09]{第 9 章}）

\textbf{Wiener filter}
维纳滤波器：根据信号和噪声的频谱来提高信噪比的最优滤波器。（见\hyperref[ch-17]{第
17 章}）

\textbf{Windowed-sinc} 加窗
sinc：用来把一个频带与另一个频带分开的数字滤波器。（见\hyperref[ch-16]{第
16 章}）

\section{Z}

\textbf{z-domain} z 域：由 z 变换定义的域。也叫 z
平面(z-plane)。（见\hyperref[ch-33]{第 33 章}）

\textbf{z-transform} z
变换：分析由差分方程支配的离散系统（例如递归（IIR）滤波器）的数学方法。把时域信号变成
z 域信号。（见\hyperref[ch-33]{第 33 章}）

\textbf{Zero} 零点：拉普拉斯变换和 z 变换中的术语。把 s 域或 z
域的传递函数写成一个多项式除以另一个多项式时，分子的根是系统的零点。另见极点(pole)。（见\hyperref[ch-32]{第
32 章}）

\textbf{Zero phase}
零相位：相位处处为零的系统。只有当冲激响应关于原点左右对称时才会出现。另见线性相位(linear
phase)。（见\hyperref[ch-07]{第 7 章}）

\textbf{Zeroth-order hold}
零阶保持：数模转换中的术语，指模拟信号在两次转换之间保持恒定值，因而呈阶梯状。（见\hyperref[ch-03]{第
3 章}）

\begin{center}\rule{0.5\linewidth}{0.5pt}\end{center}

\section{译名统一说明}

各章译文对少数术语用过不同的中文名，本表统一如下（各章正文未改动）：

{\def\LTcaptype{none} % do not increment counter
\begin{longtable}[]{@{}
  >{\raggedright\arraybackslash}p{(\linewidth - 4\tabcolsep) * \real{0.3333}}
  >{\raggedright\arraybackslash}p{(\linewidth - 4\tabcolsep) * \real{0.3333}}
  >{\raggedright\arraybackslash}p{(\linewidth - 4\tabcolsep) * \real{0.3333}}@{}}
\toprule\noalign{}
\begin{minipage}[b]{\linewidth}\raggedright
英文
\end{minipage} & \begin{minipage}[b]{\linewidth}\raggedright
本表采用
\end{minipage} & \begin{minipage}[b]{\linewidth}\raggedright
各章中出现过的其他译法
\end{minipage} \\
\midrule\noalign{}
\endhead
\bottomrule\noalign{}
\endlastfoot
circular convolution & 循环卷积 & 环形卷积（第 10 章）、圆周卷积（第 24
章） \\
recursion equation & 递归方程 & 递推方程（第 7 章） \\
deconvolution & 反卷积 & 解卷积（第 9 章） \\
edge response & 边缘响应 & 边沿响应（第 14 章） \\
Harvard architecture & 哈佛结构 & 哈佛架构（第 4 章） \\
\end{longtable}
}
